\documentclass[10pt]{beamer}
\usetheme{metropolis}
\usepackage{graphicx}
\usepackage{listings}
\usepackage{booktabs}
\usepackage{xcolor}
\definecolor{codebg}{HTML}{F5F5F5}
\definecolor{codegreen}{HTML}{2E7D32}
\definecolor{codegray}{HTML}{757575}
\definecolor{codeblue}{HTML}{1565C0}
\lstdefinestyle{rstyle}{language=R,backgroundcolor=\color{codebg},basicstyle=\ttfamily\scriptsize,
  keywordstyle=\color{codeblue}\bfseries,stringstyle=\color{codegreen},
  commentstyle=\color{codegray}\itshape,breaklines=true,frame=single,
  rulecolor=\color{codegray!30},showstringspaces=false,tabsize=2,
  morekeywords={library,read_csv,filter,select,mutate,arrange,group_by,summarise,
    count,slice_max,slice_min,pull,rename,relocate,across,where,
    drop_na,replace_na,as.numeric,as.Date,mdy,str_remove,str_trim,
    case_when,if_else,fct_reorder,fct_lump_n,fct_infreq,
    pivot_longer,pivot_wider,separate,unite,separate_rows,
    left_join,inner_join,anti_join,
    ggplot,aes,geom_col,geom_point,coord_flip,theme_minimal,labs}}
\lstdefinestyle{pythonstyle}{language=Python,backgroundcolor=\color{codebg},basicstyle=\ttfamily\scriptsize,
  keywordstyle=\color{codeblue}\bfseries,stringstyle=\color{codegreen},
  commentstyle=\color{codegray}\itshape,breaklines=true,frame=single,
  rulecolor=\color{codegray!30},showstringspaces=false,tabsize=4,
  morekeywords={import,from,as,pd,np,query,assign,sort_values,groupby,agg,
    value_counts,nlargest,dropna,fillna,to_numeric,to_datetime,
    melt,pivot_table,merge,drop_duplicates,str,replace,rename,astype}}
\title{Data Wrangling for Visualization (R)}
\subtitle{Workshop 4 --- Module 2}
\author{Prof.Asoc.\ Endri Raco}
\institute{Data Visualization for Data Scientists \\ UNYT Tirana}
\date{2025/26}
\begin{document}

\begin{frame}
  \titlepage
\end{frame}

\begin{frame}{Module Roadmap}
  \begin{enumerate}
    \item The tidyverse pipeline: import $\rightarrow$ wrangle $\rightarrow$ visualize
    \item The pipe operator: \texttt{|>} reads top-to-bottom
    \item Six dplyr verbs: filter, select, mutate, arrange, group\_by, summarise
    \item Reshaping with tidyr: pivot\_longer, pivot\_wider
    \item Common cleaning patterns for visualization
    \item Pipe-to-plot: wrangling directly into ggplot
  \end{enumerate}
  \begin{exampleblock}{Learning Outcome}
    You will wrangle any messy CSV into tidy, ggplot-ready form using
    a single piped chain of dplyr and tidyr operations.
  \end{exampleblock}
\end{frame}

\section{The tidyverse Pipeline}

\begin{frame}{Import $\rightarrow$ Wrangle $\rightarrow$ Visualize}
  \begin{center}
    \includegraphics[width=0.95\textwidth]{figures_w04m02/tidyverse_pipeline}
  \end{center}
\end{frame}

\begin{frame}{The Pipe Operator}
  \begin{center}
    \includegraphics[width=0.92\textwidth]{figures_w04m02/pipe_operator}
  \end{center}
  \begin{alertblock}{Pro-Tip}
    R 4.1+ introduced the native pipe \texttt{|>}. The tidyverse pipe
    \texttt{\%>\%} (from magrittr) still works identically. Both pass
    the left-hand result as the first argument of the right-hand
    function. Read the chain top-to-bottom: ``take apps, then filter,
    then group, then summarise, then arrange.''
  \end{alertblock}
\end{frame}

\section{Six dplyr Verbs}

\begin{frame}{The Grammar of Data Wrangling}
  \begin{center}
    \includegraphics[width=0.95\textwidth]{figures_w04m02/dplyr_verbs}
  \end{center}
\end{frame}

\begin{frame}[fragile]{filter() + select(): Subset Rows and Columns}
\begin{lstlisting}[style=rstyle]
# filter(): keep rows matching condition
apps |>
  filter(Rating > 4.0,              # AND: comma = &
         Type == "Free",
         Category %in% c("GAME", "FAMILY", "TOOLS"))

# select(): keep, drop, or rename columns
apps |>
  select(App, Category, Rating, Reviews, Type) |>  # keep 5 cols
  select(-Reviews) |>                                # drop 1 col
  rename(app_name = App)                             # rename

# Helper functions inside select():
apps |> select(starts_with("Content"),   # prefix match
               ends_with("Rating"),       # suffix match
               where(is.numeric))         # type-based
\end{lstlisting}
\end{frame}

\begin{frame}[fragile]{mutate() + case\_when(): Create and Transform}
\begin{lstlisting}[style=rstyle]
# mutate(): add or transform columns
apps |>
  mutate(
    Reviews = as.numeric(Reviews),            # type conversion
    log_reviews = log10(Reviews + 1),         # log transform
    Size_MB = str_remove(Size, "M") |>        # string cleaning
      as.numeric(),
    date = mdy(`Last Updated`),               # date parsing
    year = year(date),                        # extract component
    price_cat = case_when(                    # conditional recode
      Price == 0 ~ "Free",
      Price < 5  ~ "Under $5",
      Price < 20 ~ "$5-$20",
      TRUE       ~ "$20+"
    )
  )

# across(): apply a function to multiple columns at once
apps |> mutate(across(where(is.character), str_trim))
\end{lstlisting}
\end{frame}

\begin{frame}[fragile]{group\_by() + summarise(): Aggregate per Group}
\begin{lstlisting}[style=rstyle]
# group_by() + summarise(): one row per group
apps |>
  filter(!is.na(Rating)) |>
  group_by(Category, Type) |>           # two-variable grouping
  summarise(
    n        = n(),                      # count
    mean_r   = mean(Rating),             # mean
    median_r = median(Rating),           # median
    sd_r     = sd(Rating),               # standard deviation
    q25      = quantile(Rating, 0.25),   # 25th percentile
    .groups  = "drop"                    # ungroup after
  ) |>
  arrange(desc(mean_r))                  # sort by mean (descending)

# Quick counting shortcut
apps |> count(Category, sort = TRUE)     # = group_by + summarise(n=n())

# slice_max/min: keep top/bottom n rows
apps |> slice_max(Rating, n = 10)        # top 10 by Rating
\end{lstlisting}
\end{frame}

\section{Reshaping with tidyr}

\begin{frame}{pivot\_longer $\leftrightarrow$ pivot\_wider}
  \begin{center}
    \includegraphics[width=0.92\textwidth]{figures_w04m02/tidyr_reshape}
  \end{center}
  \begin{exampleblock}{Data Insight}
    ggplot2 requires \textbf{long (tidy)} format: each variable is a
    column, each observation is a row. If your data has years as columns
    (wide), you must \texttt{pivot\_longer()} before plotting. The
    reverse (\texttt{pivot\_wider()}) is for creating summary tables.
  \end{exampleblock}
\end{frame}

\begin{frame}[fragile]{Reshaping in Practice}
\begin{lstlisting}[style=rstyle]
# WIDE -> LONG (for ggplot)
df_wide |>
  pivot_longer(
    cols = `2020`:`2024`,        # columns to stack
    names_to = "year",            # new column for former col names
    values_to = "revenue"         # new column for values
  ) |>
  mutate(year = as.integer(year)) |>
  ggplot(aes(x = year, y = revenue, color = region)) +
  geom_line()

# LONG -> WIDE (for summary tables)
df_long |>
  pivot_wider(
    names_from = year,
    values_from = revenue
  )

# separate(): split one column into two
apps |> separate(Genres, into = c("genre1", "genre2"), sep = ";")

# separate_rows(): one row per value (explode)
netflix |> separate_rows(listed_in, sep = ", ")
\end{lstlisting}
\end{frame}

\section{Cleaning Patterns}

\begin{frame}{Six Common Cleaning Patterns}
  \begin{center}
    \includegraphics[width=0.92\textwidth]{figures_w04m02/cleaning_patterns}
  \end{center}
\end{frame}

\begin{frame}[fragile]{Cleaning the Google Play Store Dataset}
\begin{lstlisting}[style=rstyle]
apps_clean <- read_csv("googleplaystore.csv") |>
  # 1. Remove duplicates
  distinct(App, .keep_all = TRUE) |>
  # 2. Fix types
  mutate(
    Reviews  = as.numeric(Reviews),
    Installs = str_remove_all(Installs, "[+,]") |> as.numeric(),
    Size     = case_when(
      str_detect(Size, "M") ~ str_remove(Size, "M") |> as.numeric(),
      str_detect(Size, "k") ~ str_remove(Size, "k") |> as.numeric() / 1024,
      TRUE ~ NA_real_),
    Price    = str_remove(Price, "\\$") |> as.numeric(),
    date     = mdy(`Last Updated`)
  ) |>
  # 3. Filter valid
  filter(!is.na(Rating), Type %in% c("Free", "Paid")) |>
  # 4. Factor ordering
  mutate(Category = fct_lump_n(Category, 8) |> fct_infreq())
\end{lstlisting}
\end{frame}

\section{Joins}

\begin{frame}{Combining Tables with Joins}
  \begin{center}
    \includegraphics[width=0.88\textwidth]{figures_w04m02/join_types}
  \end{center}
  \begin{alertblock}{Pro-Tip}
    \texttt{left\_join()} is the default for visualization: keep all
    rows from your main table, add columns from a lookup table.
    \texttt{anti\_join()} is invaluable for finding ``what\u2019s missing''
    --- rows in table A with no match in table B.
  \end{alertblock}
\end{frame}

\section{Pipe-to-Plot}

\begin{frame}{One Chain: Raw Data $\rightarrow$ Chart}
  \begin{center}
    \includegraphics[width=0.82\textwidth]{figures_w04m02/pipe_to_plot}
  \end{center}
\end{frame}

\section{Complete Examples}

\begin{frame}[fragile]{R: Full Wrangling $\rightarrow$ ggplot Chain}
  \begin{columns}[T]
    \begin{column}{0.52\textwidth}
\begin{lstlisting}[style=rstyle]
apps_clean |>
  group_by(Category) |>
  summarise(
    mean_r = mean(Rating),
    n = n(),
    .groups = "drop") |>
  slice_max(n, n = 8) |>
  mutate(Category =
    fct_reorder(Category, mean_r),
    hl = mean_r == max(mean_r)) |>
  ggplot(aes(x = Category,
    y = mean_r, fill = hl)) +
  geom_col(width = 0.5,
    show.legend = FALSE) +
  geom_text(aes(
    label = round(mean_r, 2)),
    hjust = -0.2, size = 3,
    fontface = "bold") +
  scale_fill_manual(values =
    c("TRUE" = "#1565C0",
      "FALSE" = "#BBBBBB")) +
  coord_flip() +
  theme_minimal() +
  labs(title = "Mean Rating
    by Top Categories",
    x = NULL, y = "Mean Rating")
\end{lstlisting}
    \end{column}
    \begin{column}{0.45\textwidth}
      \includegraphics[width=\textwidth]{figures_w04m02/r_output}

      \vspace{4pt}
      {\small The entire pipeline ---
      from grouped summary to finished
      chart --- is a single piped chain.
      No intermediate variables needed.}
    \end{column}
  \end{columns}
\end{frame}

\begin{frame}[fragile]{Python Equivalent: pandas Chain}
  \begin{columns}[T]
    \begin{column}{0.52\textwidth}
\begin{lstlisting}[style=pythonstyle]
(apps_clean
    .groupby("Category")
    .agg(
        mean_r=("Rating", "mean"),
        n=("Rating", "count"))
    .nlargest(8, "n")
    .sort_values("mean_r")
    .assign(color=lambda d:
        np.where(d["mean_r"] ==
            d["mean_r"].max(),
            "#1565C0", "#BBBBBB"))
    .plot.barh(
        y="mean_r",
        color="color",  # simplified
        figsize=(6, 5),
        legend=False)
)

# Or manual for full control:
summary = (apps_clean
    .groupby("Category")
    .agg(mean_r=("Rating","mean"),
         n=("Rating","count"))
    .nlargest(8, "n")
    .sort_values("mean_r"))

ax.barh(summary.index,
    summary["mean_r"])
\end{lstlisting}
    \end{column}
    \begin{column}{0.45\textwidth}
      \includegraphics[width=\textwidth]{figures_w04m02/py_output}

      \vspace{4pt}
      {\small pandas uses method chaining
      (\texttt{.groupby().agg().nlargest()})
      as the equivalent of R\u2019s pipe.
      The \texttt{.assign()} method
      creates columns inline.}
    \end{column}
  \end{columns}
\end{frame}

\section{Synthesis}

\begin{frame}{Key Takeaways}
  \begin{enumerate}
    \item The \textbf{pipe} (\texttt{|>}) transforms nested code into
          readable top-to-bottom chains.
    \item \textbf{Six dplyr verbs} handle 90\% of wrangling: filter,
          select, mutate, arrange, group\_by, summarise.
    \item \textbf{pivot\_longer()} is required before plotting wide data
          in ggplot2; \textbf{pivot\_wider()} creates summary tables.
    \item Common cleaning: \texttt{drop\_na()}, \texttt{as.numeric()},
          \texttt{str\_remove()}, \texttt{mdy()}, \texttt{case\_when()},
          \texttt{fct\_reorder()}.
    \item \textbf{Pipe-to-plot}: wrangle and plot in a single chain
          with no intermediate variables.
    \item \texttt{left\_join()} adds context; \texttt{anti\_join()} finds
          missing data.
  \end{enumerate}
\end{frame}

\begin{frame}{Essential Readings}
  \begin{itemize}
    \item Wickham, H. \& Grolemund, G. (2023). \emph{R for Data Science},
          2nd ed., Chapters 3--5 (Data Transformation, Tidy Data).
    \item Wickham, H. (2014). ``Tidy Data.'' \emph{Journal of Statistical
          Software}, 59(10).
    \item tidyverse documentation: \texttt{https://www.tidyverse.org}
  \end{itemize}
  \vspace{6pt}
  \textbf{Next Module}: Data Wrangling for Visualization in Python (W04-M03)
\end{frame}

\end{document}
